\SC{You don't need to keep referencing the same reference after every sentence.  You can leave it to the last relevant sentence in a paragraph or the last sentence before a new reference is needed.}
%%%%%%%%%%%%%%%%%%%%%%%%%%%%%%%%%%%%%%%%%%%%%%%%%%%
\section{Introduction to Interferometers}
% ------------------------------------------------

Studying protostellar disks at millimetre wavelengths requires high-angular-resolution observations to resolve their structure. This is a very difficult process as these objects are far away, so they appear very small in the sky, and the angular resolution is limited by diffraction \citep{Monnier}. The angular resolution ($\theta_{\text{res}}$) of a single-dish radio telescope observing at wavelength $\lambda$ is given by $\theta_{\text{res}}\sim \lambda/D$, where $D$ is the diameter of the dish \citep{Monnier}. Achieving the angular resolution required to \SC{study a disk} with a single-dish telescope would require a dish kilometres in diameter, which is expensive and impractical, and it could not be steered to observe other objects. 


To overcome this limitation, radio astronomers use interferometers, where multiple widely separated dishes synthesize a large dish, improving angular resolution  \citep{Ludovici2026}. The equation above now becomes $\theta_{\text{res}}\sim \lambda/B$, where $B$ is the maximum distance between two telescopes, called the baseline \citep{Monnier}. 














%%%%%%%%%%%%%%%%%%%%%%%%%%%%%%%%%%%%%%%%%%%%%%%%%%%



 % Not only is thsi process easier in some ways as you do not need to build dishes that have a size equivalent to the length of a marathon (To obtain 1 arcsecond resolution at 21 cm wavelength, the aperature of $\sim$ 42 km (essentially a marathon) is required \citep{Perley2026}.), it is also cheaper (the FAST 500 meter disk telescope with arcminuite resolution at 11.4 GHz was 250 million CAS to build \citep{AliceCurtin}.  )
 
 
 
 
 % In astronomy, we want to know the angular distribution and the brightness distribution of the emission of a distant radiating source \citep{Perley2026}. To do this, we need to resolve the emission using an instrument of sufficient angular resolution \citep{Perley2026}. This is very difficult as our targets are very far away, the emission is weak and has a small angular size \citep{Perley2026}. 
 
 
 
% An interferometer is an indirect imaging device \citep{Rau2026}. 


% Radio telescopes sum the electric fields over an aperture of size $D$, and for this, the diffraction theory applies \citep{Perley2026}. 



% The angular resolution of a telescope is defined by:
% \begin{equation}
%     \theta_{rad} = \frac{\lambda}{D}, or \quad \theta_{arcsec} \sim \frac{2\lambda_{cm}}{D_{km}}
% \end{equation}
% \noindent where $\lambda$ is \add{ } and $D$ is \add{ } \citep{AliceCurtin} \citep{Perley2026}. One option is to build a massive single-dish telescope. (Note that these $\theta$ values are used in this equation only, and when $\theta$ symbols pop up later in this thesis, they are not the same. One easier solution is to make many smaller dishes that are spread apart, have long baselines, and large effective collections \citep{AliceCurtin}.  \question{Define dish?} \question{define effective collection area?}



% These interferometers have not only helped astronomers achieve the angular resolution they require at radio wavelengths, but have surpassed the resolution at ground- and space-based optical telescopes \citep{Monnier}. 



% One of the most popular interferometers that allows us to make high-quality images with high sensitivity and high angular resolution and a large field-of-view using an array of telescopes is the Atacama Large Millimetre Array (ALMA). \citep{Monnier}.
%%%%%%%%%%%%%%%%%%%%%%%%%%%%%%%%%%%%%%%%%%%%%%%%%%%

